%\subsection{Challenges and Techniques}
There are several criticalities that emerge when trying to combine techniques from the toolbox of regret minimization in MABs and the analysis techniques for competitive ratio.
\paragraph{Buffer Divergence between \alg and \opt} In MABs, an algorithm can always choose the same action as \opt. This is not the case in $K$-OPSD. A single different decision leads to different buffers in the subsequent rounds. Thus, it is not possible, in general, to bound the instantaneous regret of an algorithm in this setting. Bounding the $\alpha$-regret in $K$-OPSD requires a charging scheme as an intermediate step that maps a set of subsequent actions from the algorithm to a set of actions in \opt schedule. This issue is particularly critical in randomized algorithms, where the stochastic structure exhibits an interplay between the randomness of the environment and the one of the algorithm. 
\paragraph{Breaking the $\Phi$ barrier for the competitive ratio} In the $2$-bounded OPSD setting, when the number of distinct weights is finite and equal to $K$, we provide an algorithm that achieves a competitive ratio of $\theta_K < \Phi$. The values of $\theta_K$ span from $\sqrt{2}$ to $\Phi$ as $K$ grows from $2$ to $\infty$. The algorithm builds over the well-known family of $\beta$-Earliest Deadline First ($\beta$-EDF) algorithms. $\beta$-EDF uses a static threshold to decide how close to the expiration the next packet to schedule must be. Our algorithm, \algtheta, uses a dynamic threshold and a restart mechanism. The sequence of thresholds is the solution to a system of $K+1$ equations that converge to the lower bound instance from \cite{hajek2001competitiveness}. Using a dynamic threshold requires a carefully crafted charging mechanism, where the schedule of \algtheta is divided in epochs. We write the competitive ratio as the solution to a linear optimization problem, and search for the worst possible vertex of the solution space.
\paragraph{Providing a $\theta_{K-1}$-regret lower bound} We provide a lower bound for the $\theta_{K-1}$-regret in $2$-bounded $K$-OPSD instances of $\Omega\left(\sqrt{T}\right)$. Proving a lower bound on the regret in bandits usually requires two instances that are hard to distinguish (from a statistical perspective) but with different enough optimal strategies. On the other hand, proving a lower bound on the competitive ratio for OPSD requires a carefully crafted instance where every decision always leads to the same competitive ratio. The two objectives can easily collide: two instances where every decision always leads to the same outcome cannot be different enough to generate a meaningful regret lower bound. We overcome this problem by proposing a \textit{three-instance} construction where the arrival sequence is partitioned into gadgets that repeats.
\paragraph{UCBs are not enough for randomized algorithms}
Our algorithms highlight a crucial difference in the construction of deterministic learning algorithms from randomized ones. Deterministic learning algorithms only use the notion of \textit{upper confidence bound} (UCB) to estimate the expected value of a packet type. Being optimistic is enough to upper bound the $\alpha$-regret in this case. Instead, our randomized algorithms make use of both UCBs and lower confidence bounds (LCBs). Indeed, LCBs allow randomized algorithms to deal with ratios of estimates in a more natural way. This emerges from the proofs of Theorems \ref{thr:upperboundrtwo} and \ref{thr:upperboundrs}.